% GENERATED FILE. Edit canonical lessons, then run npm run generate:books.
% canonical-source-hash: fnv1a64:711bfc124d355afa
% canonical-lessons: HI-C433-qatar, HI-C433-tokan, HI-C433-kauntar, HI-C433-hat, HI-C433-thamana

\chapter{Queues and Counters}
\label{ch:hi-queues-and-counters}
\hlchaptermodality{433}

\begin{chapteropening}
\textbf{By the end of this chapter:} \emph{I can stand in a queue, take a token, wait at a counter and shop at the weekly market.}

Complete the last lesson of chapter 433: I can stand in a queue, take a token, wait at a counter and shop at the weekly market.
\end{chapteropening}

\section[\texorpdfstring{\hi{क़तार} (qatār)}{क़तार (qatār)}]{\hi{क़तार} (qatār) — a queue; a row}

\label{lesson:HI-C433-qatar}



\begin{quote}
Before the new one: say the Hindi for a label, then the Hindi for to wrap.
\end{quote}

\subsection*{You'll want to know: \hi{क़तार}}
\textbf{\hi{क़तार}} — \emph{qatār} — \textquotedblleft{}a queue; a row\textquotedblright{}.

\begin{quote}
\hi{टिकट} \hi{घर} \hi{के} \hi{सामने} \hi{लंबी} \hi{क़तार} \hi{थी।}
\end{quote}

\begin{quote}
\emph{ṭikaṭ ghar ke sāmne lambī qatār thī.}
\end{quote}

\begin{quote}
There was a long queue in front of the ticket office.
\end{quote}

Feminine.

\begin{cousinweb}
\textbf{\hi{क़तार}} (\emph{qatār}) is usually traced to Arabic \emph{qiṭār}, \textquotedblleft{}a line of camels tied one behind another\textquotedblright{}, through Persian: people waiting in line, one behind another.
\end{cousinweb}

\begin{grammarlens}[title={What you've built}]
One more word for getting around.
\end{grammarlens}

\subsection*{Your turn}
\begin{itemize}
\raggedright
  \item \emph{Say it:} \emph{qatār}
  \item \emph{Say it:} \emph{qatār}, once more
  \item \emph{Say it:} \emph{ṭikaṭ ghar ke sāmne lambī qatār thī.} — There was a long queue in front of the ticket office.
  \item \emph{Recall:} say the Hindi for a label, then the Hindi for to wrap, then say \emph{qatār} again
\end{itemize}

\begin{tcolorbox}[breakable,colback=teal!4,colframe=teal!35!black,title={Before you move on}]
What does \hi{क़तार} mean? (\textquotedblleft{}A queue; a row\textquotedblright{}.) Say it once more.
\end{tcolorbox}
\section[\texorpdfstring{\hi{टोकन} (ṭokan)}{टोकन (ṭokan)}]{\hi{टोकन} (ṭokan) — a token}

\label{lesson:HI-C433-tokan}



\begin{quote}
Before the new one: say the Hindi for to wrap, then the Hindi for a queue.
\end{quote}

\subsection*{You'll want to know: \hi{टोकन}}
\textbf{\hi{टोकन}} — \emph{ṭokan} — \textquotedblleft{}a token\textquotedblright{}.

\begin{quote}
\hi{डाकघर} \hi{में} \hi{पहले} \hi{टोकन} \hi{लीजिए।}
\end{quote}

\begin{quote}
\emph{ḍākghar meṁ pahle ṭokan lījie.}
\end{quote}

\begin{quote}
At the post office, take a token first.
\end{quote}

Masculine.

\begin{cousinweb}
\textbf{\hi{टोकन}} (\emph{ṭokan}) is borrowed from English \emph{token}, the numbered slip that keeps your place in a queue.
\end{cousinweb}

\begin{grammarlens}[title={What you've built}]
One more word for getting around.
\end{grammarlens}

\subsection*{Your turn}
\begin{itemize}
\raggedright
  \item \emph{Say it:} \emph{ṭokan}
  \item \emph{Say it:} \emph{ṭokan}, once more
  \item \emph{Say it:} \emph{ḍākghar meṁ pahle ṭokan lījie.} — At the post office, take a token first.
  \item \emph{Recall:} say the Hindi for to wrap, then the Hindi for a queue, then say \emph{ṭokan} again
\end{itemize}

\begin{tcolorbox}[breakable,colback=teal!4,colframe=teal!35!black,title={Before you move on}]
What does \hi{टोकन} mean? (\textquotedblleft{}A token\textquotedblright{}.) Say it once more.
\end{tcolorbox}
\section[\texorpdfstring{\hi{काउंटर} (kāunṭar)}{काउंटर (kāunṭar)}]{\hi{काउंटर} (kāunṭar) — a counter}

\label{lesson:HI-C433-kauntar}



\begin{quote}
Before the new one: say the Hindi for a queue, then the Hindi for a token.
\end{quote}

\subsection*{You'll want to know: \hi{काउंटर}}
\textbf{\hi{काउंटर}} — \emph{kāunṭar} — \textquotedblleft{}a counter\textquotedblright{}.

\begin{quote}
\hi{दूसरे} \hi{काउंटर} \hi{पर} \hi{जाइए।}
\end{quote}

\begin{quote}
\emph{dūsre kāunṭar par jāiye.}
\end{quote}

\begin{quote}
Please go to the second counter.
\end{quote}

Masculine.

\begin{cousinweb}
\textbf{\hi{काउंटर}} (\emph{kāunṭar}) is borrowed from English \emph{counter}, at a bank, a shop or a station.
\end{cousinweb}

\begin{grammarlens}[title={What you've built}]
One more word for getting around.
\end{grammarlens}

\subsection*{Your turn}
\begin{itemize}
\raggedright
  \item \emph{Say it:} \emph{kāunṭar}
  \item \emph{Say it:} \emph{kāunṭar}, once more
  \item \emph{Say it:} \emph{dūsre kāunṭar par jāiye.} — Please go to the second counter.
  \item \emph{Recall:} say the Hindi for a queue, then the Hindi for a token, then say \emph{kāunṭar} again
\end{itemize}

\begin{tcolorbox}[breakable,colback=teal!4,colframe=teal!35!black,title={Before you move on}]
What does \hi{काउंटर} mean? (\textquotedblleft{}A counter\textquotedblright{}.) Say it once more.
\end{tcolorbox}
\section[\texorpdfstring{\hi{हाट} (hāṭ)}{हाट (hāṭ)}]{\hi{हाट} (hāṭ) — a weekly market; a village market}

\label{lesson:HI-C433-hat}



\begin{quote}
Before the new one: say the Hindi for a token, then the Hindi for a counter.
\end{quote}

\subsection*{You'll want to know: \hi{हाट}}
\textbf{\hi{हाट}} — \emph{hāṭ} — \textquotedblleft{}a weekly market; a village market\textquotedblright{}.

\begin{quote}
\hi{हर} \hi{मंगलवार} \hi{को} \hi{गाँव} \hi{में} \hi{हाट} \hi{लगती} \hi{है।}
\end{quote}

\begin{quote}
\emph{har maṅgalvār ko gā̃v meṁ hāṭ lagtī hai.}
\end{quote}

\begin{quote}
Every Tuesday there is a market in the village.
\end{quote}

Feminine.

\begin{cousinweb}
\textbf{\hi{हाट}} (\emph{hāṭ}) is a tadbhava of Sanskrit \emph{haṭṭa}, \textquotedblleft{}a market, a fair\textquotedblright{} (CDIAL 13944), which lists Hindi \emph{hāṭ}, \textquotedblleft{}market, shop\textquotedblright{}; \hi{हड़ताल}, a strike, closing the market, is in the same entry.
\end{cousinweb}

\begin{grammarlens}[title={What you've built}]
One more word for getting around.
\end{grammarlens}

\subsection*{Your turn}
\begin{itemize}
\raggedright
  \item \emph{Say it:} \emph{hāṭ}
  \item \emph{Say it:} \emph{hāṭ}, once more
  \item \emph{Say it:} \emph{har maṅgalvār ko gā̃v meṁ hāṭ lagtī hai.} — Every Tuesday there is a market in the village.
  \item \emph{Recall:} say the Hindi for a token, then the Hindi for a counter, then say \emph{hāṭ} again
\end{itemize}

\begin{tcolorbox}[breakable,colback=teal!4,colframe=teal!35!black,title={Before you move on}]
What does \hi{हाट} mean? (\textquotedblleft{}A weekly market; a village market\textquotedblright{}.) Say it once more.
\end{tcolorbox}
\section[\texorpdfstring{\hi{थमाना} (thamānā)}{थमाना (thamānā)}]{\hi{थमाना} (thamānā) — to hand (to); to put into someone's hand}

\label{lesson:HI-C433-thamana}



\begin{quote}
Before the new one: say the Hindi for a counter, then the Hindi for a weekly market.
\end{quote}

\subsection*{You'll want to know: \hi{थमाना}}
\textbf{\hi{थमाना}} — \emph{thamānā} — \textquotedblleft{}to hand (to); to put into someone's hand\textquotedblright{}.

\begin{quote}
\hi{दुकानदार} \hi{ने} \hi{मुझे} \hi{थैला} \hi{थमाया।}
\end{quote}

\begin{quote}
\emph{dukāndār ne mujhe thailā thamāyā.}
\end{quote}

\begin{quote}
The shopkeeper handed me the bag.
\end{quote}

Past, with \hi{ने} (he or she …): \hi{उसने} \hi{मुझे} \hi{चाबी} \hi{थमाई}, \emph{usne mujhe cābī thamāī}.

\begin{cousinweb}
\textbf{\hi{थमाना}} (\emph{thamānā}) is the causative of \hi{थामना}, \textquotedblleft{}to hold\textquotedblright{}, which CDIAL lists under Sanskrit \emph{stambhate}, \textquotedblleft{}supports, arrests\textquotedblright{} (CDIAL 13683): to make someone hold a thing.
\end{cousinweb}

\begin{grammarlens}[title={What you've built}]
One more word for getting around.
\end{grammarlens}

\subsection*{Your turn}
\begin{itemize}
\raggedright
  \item \emph{Say it:} \emph{thamānā}
  \item \emph{Say it:} \emph{thamānā}, once more
  \item \emph{Say it:} \emph{dukāndār ne mujhe thailā thamāyā.} — The shopkeeper handed me the bag.
  \item \emph{Recall:} say the Hindi for a counter, then the Hindi for a weekly market, then say \emph{thamānā} again
\end{itemize}

\begin{tcolorbox}[breakable,colback=teal!4,colframe=teal!35!black,title={Before you move on}]
What does \hi{थमाना} mean? (\textquotedblleft{}To hand (to); to put into someone's hand\textquotedblright{}.) Say it once more.
\end{tcolorbox}
